\documentclass[11pt]{article}
\usepackage[no-math]{fontspec}
\usepackage[french,english]{babel}
\usepackage{amsmath}
\usepackage{amssymb}
\usepackage{amsthm}
\usepackage{mathrsfs}
\usepackage{tikz-cd}
\usepackage{enumitem}
\usetikzlibrary{babel}

\newcommand{\printedpage}[1]{\marginpar{\footnotesize p.~#1}}
\newcommand{\msp}[1]{\marginpar{\footnotesize [#1]}}
\newcommand{\HOT}{\mathrm{HOT}}
\newcommand{\Cat}{\mathrm{Cat}}
\newcommand{\Ob}{\mathrm{Ob}}
\newcommand{\Fl}{\mathrm{Fl}}
\newcommand{\id}{\mathrm{id}}
\newcommand{\cA}{\mathcal{A}}
\newcommand{\ul}[1]{\underline{\mathrm{#1}}}
\newcommand{\circled}[1]{\tikz[baseline=(c.base)]\node[draw,circle,inner sep=1pt](c){#1};}

\begin{document}

\printedpage{39}
% head: A. Grothendieck : Dérivateurs, Ch. VI, Édition des Universités de Montpellier II et Paris VII

\[
\begin{tikzcd}[column sep=large, row sep=large]
 & \HOT_W(S)\approx(\Cat/S)(W_S^{\mathrm{d}})^{-1} \\
\HOT^{\mathrm{lc}}_W(S)\approx(\Cat/S)W_S^{-1} \arrow[ur, "i"] \arrow[r, "C"] \arrow[dr, "j"'] & (\ul{Parf}\,S)W_S^{-1} \arrow[u, "\text{\footnotesize since $W_S=W_S^{\mathrm{d}}$ on $\ul{Parf}$}"'] \arrow[d, "\text{\footnotesize since $W_S=W_S^{\mathrm{g}}$ on $\ul{Parf}$}"] \\
 & \HOT^{o}_W(S)\approx(\Cat/S)(W_S^{\mathrm{g}})^{-1}.
\end{tikzcd}
\]

\msp{page 64}
We see from this that we have canonical isomorphisms
\[
\rho i\simeq\id,\qquad\sigma j\simeq\id
\]
in the category of endofunctors of $(\Cat/S)W_S^{-1}$ ($\approx\HOT^{\mathrm{lc}}_W(S)$), an isomorphism which is made explicit here by the fact that $\rho i$, $\sigma j$ are deduced from the same functor $C$ by passage to the localized categories $W_S^{-1}$ (so in fact $\rho i=\sigma j$ !), and that we have an arrow functorial in $X$
\[
X\xrightarrow{i_X}C(X),
\]
which is in $W$ for every $X$. Thus, $\rho$ and $\sigma$ are, up to isomorphism, \textit{retractions} of the categories $\HOT_W(S)$, $\HOT^{o}_W(S)$ onto the category $\HOT^{\mathrm{lc}}_W(S)$. Moreover, these retractions are what are called localization functors (as we see from the diagram on page 62). (We say that $C\xrightarrow{F}C'$ is a \textit{localization functor} if, denoting by $\Sigma\subseteq\Fl\,C$ the set of arrows of $C$ transformed into isomorphisms by $F$, the functor
\[
C\Sigma^{-1}\xrightarrow{\bar{F}}C'
\]
which factors $F$ is an equivalence of categories.)

\msp{page 65}
We know, moreover, that if a functor $f$ has an adjoint $g$ (left or right), then $f$ is a functor \texttt{[of]} localization if and only if the adjoint $g$ is fully faithful. To prove that $i$, $j$ are fully faithful, it would therefore be tempting to prove that they are adjoint to the functors $\rho$, $\sigma$.

To get our bearings, I shall set
\[
\left\{
\begin{aligned}
&\cA=\Cat/S,\ \ \Sigma_0=W_S^{\mathrm{d}},\ \ \Sigma=W_S\qquad\quad(\subseteq\Fl(\cA))\\
&C:\cA\longrightarrow\cA\ \text{ the familiar functor, which has the virtue}\\
&\qquad\qquad\qquad\quad C(\Sigma)\subseteq\Sigma_0\quad.
\end{aligned}
\right.
\]

We thus have the situation
\printedpage{40}
% head: A. Grothendieck : Dérivateurs, Ch. VI, Édition des Universités de Montpellier II et Paris VII
\[
(*)\qquad\qquad(\cA,\Sigma_0)\ \underset{C}{\overset{\rho_0\,=\,\id_{\cA}}{\rightleftarrows}}\ (\cA,\Sigma),
\]
inducing by passage to the localized categories
\[
\cA\Sigma_0^{-1}\quad\underset{\bar{C}}{\overset{\rho\,=\,\bar{\rho}_0}{\rightleftarrows}}\ \cA\Sigma^{-1},
\]
we want to check that $\rho$ and $C$ \texttt{[rather $\bar{C}$]} are adjoint to one another, $C$ \texttt{[rather $\bar{C}$]} \textit{surely on the right,} which means that we must have functorial homomorphisms
\[
\left\{
\begin{aligned}
&\id_{\cA\Sigma^{-1}}&&\rightleftarrows&&\rho\,\bar{C}\quad(=\overline{\rho_0C})\\
&\id_{\cA\Sigma_0^{-1}}&&\rightleftarrows&&\bar{C}\rho\quad(=\overline{C\rho_0})
\end{aligned}
\right.
\]
(\footnote[41]{but if $C$ is on the \textit{right,} the arrows here must go in the opposite direction and this is indeed what I must resign myself to on p.~67.}) \texttt{[the arrows at the bottom go in the corrected direction, but the arrows at the top are not erased]} satisfying two compatibility relations to which we shall return. We may hope to find these functorial homomorphisms starting from
\[
\left\{
\begin{aligned}
&\id_{\cA}&&\xrightarrow{\ i\ }&&\rho_0C&&=&&C\\
&\id_{\cA}&&\xleftarrow{\ p\ }&&C\rho_0&&=&&C
\end{aligned}
\right.
\]

\msp{page 66}
and in our case we do indeed have
\[
i:\id_{\cA}\longrightarrow C\qquad\text{i.e. }\ i_X:X\longrightarrow C(X)
\]
functorial in $X\in\Ob\,\cA$ (and moreover $i_X$ is in $\Sigma$). There is however no natural arrow in the opposite direction $p_X$
\[
C(X)\xrightarrow{\ ?\ }X.
\]
I had thought of the projection $\pi_X:C(X)\longrightarrow X$, deduced from the construction $C(X)\simeq X\times_S(\ul{Ch}_\infty(S),\sigma)$, but it is unfortunately not an $S$-morphism. As we are looking for an arrow in $\cA\Sigma_0^{-1}$, we could make do with having a functorial arrow in the opposite direction, provided it is in $\Sigma_0$, hence can be inverted when we work in $\cA\Sigma_0^{-1}$. The obvious candidate is $i_X:X\longrightarrow C(X)$ which we have just spoken of. The trouble is that $i_X$ is only in $\Sigma$, not in $\Sigma_0$ (\footnote[42]{If this were not so, the relative homotopy type over $S$ defined by $X/S$ would be locally constant (since this is so for the one defined by $C(X)$), which is not always the case.}). Yet, when $X$ is itself of the form $C(Y)$ (more generally, if it is Hot-fibered over $S$), then the fact that $X\longrightarrow C(X)$ is in $\Sigma$ already implies that it is in $\Sigma_0$. Let us therefore note this % [sic] evident -> évident, accent missing
\[
\left\{
\begin{array}{lrl}
i:\id_{\cA}\longrightarrow C\ \text{such that} & 1^o) & i_X:X\longrightarrow C(X)\ \textit{in}\ \Sigma\\[1ex]
 & & \text{for every }X\in\Ob\,\cA\\[1ex]
 & 2^o) & i_X\ \text{in }\Sigma_0\text{ if }X\text{ of the}\\[1ex]
 & & \text{form }C(Y)
\end{array}
\right.
\]

\printedpage{41}
% head: A. Grothendieck : Dérivateurs, Ch. VI, Édition des Universités de Montpellier II et Paris VII

\msp{page 67}
The idea comes to me, after some shitty attempts \texttt{[sic]}, that $\bar{C}$ might be the right adjoint, and not the left \texttt{[this is the corrected version (right and left interchanged) cf.~p.~65]}, i.e. that the arrows sought might go in the opposite direction
\[
\left\{
\begin{aligned}
&\id_{\cA\Sigma^{-1}}&&\xleftarrow{\ p\ }&&\rho\,\bar{C}\quad(=\overline{\rho_0C})\\
&\id_{\cA\Sigma_0^{-1}}&&\xrightarrow{\ i\ }&&\bar{C}\rho\quad(=\overline{C\rho_0})
\end{aligned}
\right.
\]
(so I reverse the notations for $p$ and $i$), thus we are led if possible to look for arrows
\[
\begin{aligned}
&\id_{\cA}&&\xleftarrow{\ p\ }&&\rho_0C&&=&&C\\
&\id_{\cA}&&\xrightarrow{\ i\ }&&C\rho_0&&=&&C,
\end{aligned}
\]
thus arrows functorial in $X$
\[
X\xrightarrow{i_X}C(X)\xrightarrow{p_X}X.
\]
We do indeed have a functorial $i_X$ - always the same one! For $p_X$, we have no more than before, but recalling that it suffices to find it in $\cA\Sigma^{-1}$ (and no longer $\cA\Sigma_0^{-1}$!), it suffices to invert the arrow $i_X$, which is indeed an \textit{isomorphism} in $\cA\Sigma^{-1}$ (but not in $\cA\Sigma_0^{-1}$). We shall therefore set
\[
p_X=i_X^{-1}:C(X)\longrightarrow X\qquad\quad\text{in }\cA\Sigma^{-1}.
\]
We have to prove just two compatibilities, for the functors $\bar{C}$ and $\rho$ between $\cA\Sigma_0^{-1}$ and $\cA\Sigma^{-1}$ (and I shall write $C$ instead of $\bar{C}$): the two composites
\[
\left\{
\begin{array}{l}
\begin{tikzcd}[ampersand replacement=\&]
\rho \arrow[r, "\rho*i"] \arrow[rr, bend left=20, dotted, "\id"] \& \rho\,C\rho \arrow[r, "p*\rho"] \& \rho
\end{tikzcd}
\qquad\quad\text{\footnotesize (relation for functors with values in $\cA\Sigma^{-1}$)}\\[2ex]
\begin{tikzcd}[ampersand replacement=\&]
C \arrow[r, "i*C"] \arrow[rr, bend right=20, dotted, "\id"'] \& C\rho\,C \arrow[r, "C*p"] \& C
\end{tikzcd}
\qquad\quad\text{\footnotesize (relation for functors with values in $\cA\Sigma_0^{-1}$)}
\end{array}
\right.
\]
must be the identity $\id_\rho$ resp. $\id_C$.

\msp{page 68}
The first relation is checked tautologically, it amounts to saying that in $\cA\Sigma^{-1}$, we have an identity composite of the
\[
X\xrightarrow{i_X}C(X)\xrightarrow{i_X^{-1}}X\qquad\quad!
\]
The second means that for every $X$, we have an identity composite in $\cA\Sigma_0^{-1}$
\[
C(X)\xrightarrow{i_{C(X)}}C(C(X))\xrightarrow{C(p_X)}C(X).
\]
Now $p_X:C(X)\longrightarrow X$ had been defined in $A\Sigma^{-1}$ as the inverse $i_X^{-1}$ of $i_X$, on the other hand the functor $C$ on $\cA$ defines by hypothesis $\cA\Sigma^{-1}\longrightarrow\cA\Sigma_0^{-1}$, more precisely $C(\Sigma)\subseteq\Sigma_0$ % [sic] $A\Sigma^{-1}$ (italic A) -> $\mathcal{A}\Sigma^{-1}$, elsewhere script A
\printedpage{42}
% head: A. Grothendieck : Dérivateurs, Ch. VI, Édition des Universités de Montpellier II et Paris VII
thus $C(i_X)$ is invertible in $\cA\Sigma_0^{-1}$, and $C(p_X)=C(i_X)^{-1}$. So the compatibility amounts to the relation
\[
C(i_X)^{-1}i_{C(X)}=\id\qquad\quad\textit{in}\ \ \cA\Sigma_0^{-1},
\]
i.e.
\[
\boxed{\,i_{C(X)}=C(i_X)\,}:C(X)\longrightarrow C^2(X).
\]
In other words, we must prove the: % [sic] d'autre termes -> d'autres termes, plural

\textbf{Lemma.} \textit{Let $X$ be in $\Cat/S$, let us consider the two arrows}
\[
i_{C(X)},\ C(i_X):C(X)\rightrightarrows C(C(X)).
\]
\textit{These arrows are equal, if not in the absolute sense, at least in $(\Cat/S)(W_S^{\mathrm{d}})^{-1}$.}

Here $f_X:X\longrightarrow S$ is the standard projection and
\[
\begin{aligned}
C(X)\ &=\ \quad\{(x,c)\ \mid\ x\in\Ob\,X,\ c:f_X(x)\ \underbrace{\longrightarrow}_{\text{path in }S}\ t\}\\
C^2(X)\ &=\ \{(x,c,c')\ \mid\ x\in\Ob\,X,\ \underbrace{f_X(x)\xrightarrow{c}t\xrightarrow{c'}t'}_{c,\ c'\text{ paths in }S}\},
\end{aligned}
\]
the projections to $S$ are given by
\[
\begin{aligned}
(x,c)&&\longmapsto&&t&&=&&\text{target }c\\
(x,c,c')&&\longmapsto&&t'&&=&&\text{target }c'
\end{aligned}
\]

\msp{page 69}
We find
\[
i_{C(X)}(x,c)=(x,c,1_t):
\begin{tikzcd}
x \arrow[d] & & \\
f_X(x) \arrow[r, "c"] & t \arrow[r, "1_t", "\text{empty path}"'] & t
\end{tikzcd}
\]
\[
C(i_X)(x,c)=(x,1_{f_X(x)},c):
\begin{tikzcd}
x \arrow[d] & & \\
f_X(x) \arrow[r, "1_{f_X(x)}", "\text{empty path}"'] & f_X(x) \arrow[r, "c"] & t
\end{tikzcd}
\]
But we have a choice between two theories of paths, to construct $C(X)$: either with the $\ul{Ch}_\infty(S)$, or with the $\ul{Ch}(S)$. ``It gives the same results'', but depending on the situation one is more convenient than the other. I shall work with the $\ul{Ch}(S)$, in which case the $1_S$ are left and right units, and we have a composition law for paths. This law allows us to define
\[
\begin{aligned}
C^2(X)&&\xrightarrow{\ \gamma\ }&&C(X)\\
(x,c,c')&&\longmapsto&&(x,c'\circ c),
\end{aligned}
\]
and the composites of $i_{C(X)}$ and of $C_{i(X)}$ with this arrow are $\id_{C(X)}$ in both cases % [sic] $C_{i(X)}$ -> $C(i_X)$, the arrow considered is C(i_X)
\printedpage{43}
% head: A. Grothendieck : Dérivateurs, Ch. VI, Édition des Universités de Montpellier II et Paris VII
\[
\begin{tikzcd}
C(X) \arrow[r, shift left] \arrow[r, shift right] \arrow[rr, bend right=15, "\id"'] & C^2(X) \arrow[r, "\gamma"] & C(X).
\end{tikzcd}
\]

So to show that $i_{C(X)}=C(i_X)$ in $(\Cat/S)(W_S^{\mathrm{d}})^{-1}$, it suffices to prove that $\gamma$ is in $W_S^{\mathrm{d}}$. But we know that the functor C transforms $W_S$ into $W_S^{\mathrm{d}}$, so $C(i_X)$ is in $W_S^{\mathrm{d}}$, and since $\gamma\circ C(i_X)=\id$, $\gamma$ is in $W_S^{\mathrm{d}}$, OK! % [sic] foncteur C (upright roman C) -> foncteur $C$, math italic elsewhere

\msp{page 70}
I feel like grouping what has been proved into a summary statement.

\textbf{Theorem-Scholie.} We consider the diagram of categories on page 8 (\footnote[43]{We assume for $W$ the conditions a) b) c) of p.~45, in addition L5 (p.~43). Cf. comments p.~81.}).

\circled{I} When we equip all 13 categories of the diagram with the localizer $W_S$ (coarse $W$-equivalence, i.e. $W$-equivalence in Cat, nothing more), \textit{all} the functors of the diagram induce equivalences on the localized categories (\footnote[44]{Moreover, if $\mathcal{F}$ is a full subcategory of $\Cat/S$ which contains one of the two subcategories $\ul{Fib}_0S$, $\ul{Cofib}_0S$ (which includes the case of the 8 categories of the diagrams, other than $\Cat/S$ itself), the inclusion $(\mathcal{F},W_S|\mathcal{F})$ into $(\Cat/S,W_S)$ induces an equivalence for the localized categories.}). Let us denote by
\[
\HOT^{\mathrm{lc}}_W(S)
\]
``the'' common localized category.

\circled{II} Let us equip the four left bordering categories
\[
(*)\qquad\qquad
\begin{tikzcd}
\Cat/S & \ul{Liss}\,S \arrow[l, hook'] & \ul{Fib}\,S \arrow[l, hook'] & \ul{Fibsc}\,S \arrow[l, hook']
\end{tikzcd}
\]
with the localizer $W_S^{\mathrm{d}}$ of \textit{right} fine $W$-equivalence over $S$ ( $\begin{tikzcd}[row sep=small, column sep=small]
X \arrow[rr, "f"] \arrow[dr] & & Y \arrow[dl]\\
 & S &
\end{tikzcd}$ is in $W_S^{\mathrm{d}}$ if and only if the $s\backslash X\longrightarrow s\backslash Y$ are \texttt{[rather, iff they are in $W$]}, for every $s\in\Ob\,S$). The functors of the diagram, i.e. those of $(*)$, between these categories, hence also their composites, induce equivalences between the localized categories (\footnote[45]{More generally, this works for every full subcategory of $\Cat/S$ which contains $\ul{Fib}\,S$ (and dually, with $W_S^{\mathrm{g}}$ instead of $W_S^{\mathrm{d}}$, if $\mathcal{F}$ contains $\ul{Cofib}\,S$).}). Let us denote by
\[
\HOT_W(S)
\]
``the'' common localized category.

Dual statement for
\[
(*')\qquad\qquad
\begin{tikzcd}
\Cat/S & \ul{Prop}\,S \arrow[l, hook'] & \ul{Cofib}\,S \arrow[l, hook'] & \ul{Cofibsc}\,S, \arrow[l, hook']
\end{tikzcd}
\]

% foot (each page): Version 9 novembre 2001

\end{document}
